% \begin{document}
\chapter{Classificazione delle Rappresentazioni Irriducibili di $S_n$}
L'idea e' quella di usare polinomi simmetrici (modo vecchio) e \textit{tableau} (piu' moderno) per poter descrivere le rappresentazioni irriducibili di $ S_n $ (e quindi il loro carattere) $ \forall n \in \mathbb{N} $.

\section{Partizioni e Ordini}
\dfn{Partizione}{
  Diciamo che $ \lambda $ e' una \textit{partizione} di $ n \in \mathbb{N} $ e lo indichiamo con $ \lambda \vdash n $ se:
  \[
    \lambda = (\lambda_1, ..., \lambda_k)
  \]
  dove
  \[
  \lambda_1 \geq ... \geq \lambda_k \geq 0 \in \mathbb{N}\ \land\ \sum_{i = 1}^{k} \lambda_i = n
  \]

  Indichiamo con $ Par(n)=\{\lambda | \lambda \vdash n\} $ e con $ Par = \bigcup_{n \in \mathbb{N}} Par(n) $.\\

  Inoltre, $ |\lambda| = n $ e $ l(\lambda) = k $ e' la lunghezza della partizione.
}
L'ordinamento sui lambda nella descrizione della partizione serve a fare in modo che non vengano contate partizioni uguali ma scritte in modo diverso.
\ex{Partizione}{
  $(3,2,2) \vdash 7$
}

Indichiamo adesso tre relazioni d'ordine sulle partizioni:

\dfn{Partizione "sghemba"}{
  La relazione totale $ \mu \subseteq \lambda $ su $ Par $ e' definita come
  \[
  \forall i. \mu_i \leq \lambda_i
  \]
  Se $ l(\mu) \neq l(\lambda) $ basta estendere la partizione piu' corta con degli zeri.\\

  In questo caso, la \textit{partizione sghemba} (che non e' una vera partizione) e'
  \[
    \lambda/\mu \coloneq (\lambda_1-\mu_1, ..., \lambda_k - \mu_k)
  \]
}
\nt{
  Si puo' vedere $ \mu \subseteq \lambda $ come il fatto che il \textit{diagramma di Ferrer} di $ \mu $ e' contenuto dentro a quello di $ \lambda $. La partizione sghemba e' quindi la differenza dei due diagrammi.
}

\ex{Partizione sghemba}{
  $ \lambda = (3,2,2), \mu = (2) $, quindi abbiamo che $ \mu \subseteq \lambda $. I due diagrammi di Ferrer sono:
  \begin{center} 
    \ydiagram{2,2,3}
    \ydiagram{2}
  \end{center}
  Possiamo quindi calcolare $ \lambda / \mu $:
  \begin{center}
    \ydiagram{2,2,2+1}
  \end{center}
}

\dfn{Ordine dominante}{
  Relazione d'ordine su $ Par(n) $ indicata con $ \mu \leq \lambda $ ed e' definita come
  \[
  \forall i.\ \mu_1 + ... + \mu_i \leq \lambda_1 + ... + \lambda_i
  \]
  E' una relazione parziale in quanto se $ l(\lambda) \neq l(\mu) $ allora le due partizioni di $ n $ sono incompatibili.
}
\ex{Ordine dominante}{
  $(2,2) <= (3,1)$
}

\dfn{RLO}{
  Il Reverse Lexicografic Order e' un ordine totale su $ Par(n) $ indicato con $ \mu \leq^R \lambda $ se:
  \begin{itemize}
  \item $ \mu = \lambda $, oppure
  \item $ \exists i.\ \mu_1 = \lambda_1,...,\mu_i = \lambda_i $ ma $ \mu_{i+1} < \lambda_{i+1} $
  \end{itemize}
}
\ex{RLO su $ Par(5) $}{
  $ (5) >^R (4,1) >^R (3,2) >^R (3,1,1) >^R (2,2,1) >^R (2,1^3) >^R (1^5) $
}

\mprop{}{
  Se $ \lambda \leq \mu $, allora $ \lambda \leq^R \mu $.
}

\section{Tableau di Young}

\dfn{Semi-Standard Young Tableau (SSYT)}{
  Dati $\mu \subseteq \lambda$ partizioni, un \textbf{Tableau di Young Semistandard} di forma
  $\lambda/\mu$ è un riempimento del diagramma di Ferrers di $\lambda/\mu$ con interi positivi
  tali che le entrate siano:
  \begin{itemize}
    \item \textbf{debolmente crescenti} lungo ogni riga (da sinistra verso destra);
    \item \textbf{strettamente crescenti} lungo ogni colonna (dal basso verso l'alto,
          notazione francese).
  \end{itemize}
}

\dfn{Tableau iniettivo}{
Data una partizione $\lambda \vdash n$ con diagramma di Ferrers associato, un \textbf{tableau iniettivo} $T$ di forma $\lambda$ è un riempimento delle celle del diagramma con gli interi $1, \dots, n$, ciascuno usato esattamente una volta.
}

\ex{Tableau Iniettivo}{
  $ \lambda = (3,2) \vdash 5 $
  \begin{center}
  \begin{ytableau}
    2 & 3\\
    5 & 1 & 4
  \end{ytableau}
  \end{center}
}

\dfn{Tableau di Young Standard (SYT)}{
  Un \textbf{Tableau di Young Standard} di forma $\lambda/\mu$ con $|\lambda/\mu|=n$ è un SSYT
  in cui ogni intero da $1$ a $n$ compare \textbf{esattamente una volta}. Equivalentemente,
  è un tableau iniettivo con entrate strettamente crescenti lungo ogni riga
  (da sinistra verso destra) e lungo ogni colonna (dal basso verso l'alto).
}

\ex{Tableaux}{
  Prendiamo $ \lambda = (3,2,1) $ e $ \mu = (1) $ partizioni. Consideriamo il diagramma di Ferrer della loro partizione sghemba $ \lambda/\mu $:
  \begin{center}
    \ydiagram{1,2,1+2}
  \end{center}
  Vediamo ora come possiamo riempirlo:
  \begin{center}
  \begin{ytableau}
    1\\
    1 & 2\\
    \none & 3 & 4
  \end{ytableau}\quad\quad
  \begin{ytableau}
    4\\
    1 & 2\\
    \none & 1 & 1
  \end{ytableau}\quad\quad
  \begin{ytableau}
    2\\
    1 & 4\\
    \none & 3 & 5
  \end{ytableau}\quad\quad
  \end{center}
  Il primo non e' neanche un SSYT, il secondo lo e' ma non un SYT, il terzo invece e' proprio un SYT.\\

  Possiamo associare al secondo tableau il seguente monomio:
  \[
  X_1^3X_2^2X_3^0X_4^1X_5^0
  \]
}

L'azione naturale di $S_n$ sui tableau iniettivi di forma $\lambda$ è data da: per ogni $\sigma \in S_n$, il tableau $\sigma \cdot T$ si ottiene sostituendo ogni entrata $i$ di $T$ con $\sigma(i)$.

\ex{Azione naturale di $ S_n $ sui tableau}{
  \begin{center}
    $ (1,3) $
    \begin{ytableau}
    1 \\
    2 & 3
    \end{ytableau} = 
    \begin{ytableau}
    3\\
    2 & 1
    \end{ytableau}
  \end{center}
}

\dfn{Gruppo di riga e gruppo di colonna}{
Sia $T$ un tableau iniettivo di forma $\lambda = (\lambda_1, \dots, \lambda_k) \vdash n$.
\begin{itemize}
\item Il \textbf{gruppo di riga} $R(T) \leq S_n$ è il sottogruppo delle permutazioni che spostano le entrate di $T$ solo all'interno delle rispettive righe. Si ha $R(T) \cong S_{\lambda} := S_{\lambda_1} \times \cdots \times S_{\lambda_k}$.
\item Il \textbf{gruppo di colonna} $C(T) \leq S_n$ è il sottogruppo delle permutazioni che spostano le entrate di $T$ solo all'interno delle rispettive colonne. Si ha $C(T) \cong S_{\lambda'}$, dove $\lambda'$ è la partizione coniugata.
\end{itemize}
}

\ex{Gruppi di riga e colonna}{
  Dato $ T = $
  \begin{ytableau}
    4 & 5\\
    1&2&3
  \end{ytableau}, abbiamo che:
  \[
    R(T) = \{g \in S_5 | g: \{123\} \to \{123\} \lor \{45\} \to \{45\}\} \cong S_3 \times S_2  
  \]
  e 
  \[
    C(T) = \{g \in S_n | g: \{41\} \to \{41\} \lor \{52\} \to \{52\}\} \cong S_2 \times S_2
  \]
}

\subsection{Proprietà di coniugio}

\mprop{Coniugio dei gruppi di riga e colonna}{
Per ogni $\sigma \in S_n$ e ogni tableau iniettivo $T$:
\[
R(\sigma \cdot T) = \sigma\, R(T)\, \sigma\inv, \qquad C(\sigma \cdot T) = \sigma\, C(T)\, \sigma\inv.
\]
}

\pf{Dimostrazione}{
Dimostriamo $R(\sigma \cdot T) = \sigma R(T) \sigma\inv$; il caso $C(T)$ è identico, sostituendo ``riga'' con ``colonna''.

\begin{enumerate}
\item ($\supseteq$) Sia $p \in R(T)$. Dato $x = \sigma(y)$ nella riga di $\sigma \cdot T$, si ha
\[
(\sigma p \sigma\inv)(x) = \sigma(p(y)).
\]
Poiché $p(y)$ sta nella stessa riga di $y$ in $T$, la sua immagine $\sigma(p(y))$ sta nella stessa riga di $x$ in $\sigma \cdot T$. Dunque $\sigma p \sigma\inv \in R(\sigma \cdot T)$.

\item ($\subseteq$) Sia $q \in R(\sigma \cdot T)$. Poniamo $p := \sigma\inv q \sigma$. Per ogni $y$ in una riga di $T$, l'elemento $\sigma(y)$ sta nella riga corrispondente di $\sigma \cdot T$, quindi $q(\sigma(y))$ vi rimane, e $\sigma\inv(q(\sigma(y))) = p(y)$ torna nella riga di $y$ in $T$. Dunque $p \in R(T)$ e $q = \sigma p \sigma\inv$.
\end{enumerate}
}

\section{Lemma di partizione dei tableau}

\mlenma{Partizione dei tableau}{
Siano $T, \wt{T}$ tableau iniettivi di forma $\lambda, \wt{\lambda} \vdash n$, con $\lambda \not>^R \wt{\lambda}$. Allora si verifica esattamente una delle seguenti:
\begin{enumerate}
\item Esistono $i \neq j$ nella stessa colonna di $T$ e nella stessa riga di $\wt{T}$.
\item $\lambda = \wt{\lambda}$ e esistono $\wt{p} \in R(\wt{T})$, $q \in C(T)$ tali che $\wt{p}\wt{T} = qT$.
\end{enumerate}
}

\ex{Partizione tableau caso 1}{
  Prendiamo $ \lambda = \wt{\lambda} = (2,1) \vdash 3 $ e 
  \begin{center}
    $ T = $
  \begin{ytableau}
    2\\
    1&3
  \end{ytableau}, $ \wt{T} =  $
    \begin{ytableau}
      3\\
      1&2
    \end{ytableau}
  \end{center}

  In questo caso abbiamo che 1 e 2 sono nella stessa colonna in $ T $ e nella stessa riga in $ \wt{T} $, quindi siamo nel primo caso. Ma sappiamo anche che $ \lambda = \wt{\lambda} $ quindi non possono esistere permutazioni $ \wt{p}, q $ tali che:
  \[
    \wt{p}\wt{T} = qT
  \]
  altrimenti varrebbe anche il secondo.
}

\ex{Partizioni tableau caso 2}{
  Prendiamo $ \lambda = \wt{\lambda} = (3,2,1) \vdash 6 $ e 
  \begin{center}
    $ T = $
  \begin{ytableau}
    3\\
    1&4\\
    2&6&5
  \end{ytableau}, $ \wt{T} =  $
    \begin{ytableau}
      2\\
      6&3\\
      4&5&1
    \end{ytableau}
  \end{center}
  Siamo nel secondo caso, quindi sappiamo che $ \exists \wt{p} \in R(\wt{T}), q \in C(T) $ tali che $ \wt{p}\wt{T} = qT $, che possiamo trovare seguendo un semplice algoritmo:

  \begin{enumerate}
    \item Trovo le permutazioni di colonna per avere gli stessi numeri di $ \wt{T} $ sulla prima riga. Scelgo $ q_1 = (12)(46) $:
     \begin{center}
$ q_1T = $
  \begin{ytableau}
    3\\
    2&6\\
    1&4&5
  \end{ytableau} 
      \end{center}
    \item Ripeto per tutte le righe sopra (tranne l'ultima che sara' automatica). Basta scegliere $ q_2 = (23) $:
     \begin{center}
$ q_2q_1T = $
  \begin{ytableau}
    2\\
    3&6\\
    1&4&5
  \end{ytableau} 
      \end{center}
      Quindi $ q = q_2q_1 = (123)(46) $
    \item Per ogni riga di $ \wt{T} $ trovo le permutazioni di riga per mettere i numeri nello stesso ordine di $ qT $. Scelgo $ \wt{p} = (36)(154) $:
      \begin{center}
        $ \wt{p}\wt{T} = $
    \begin{ytableau}
      2\\
      3&6\\
      1&4&5
    \end{ytableau} $ = qT $
      \end{center}
  \end{enumerate}
}

\pf{Dimostrazione}{
Assumiamo che (1) sia falsa e mostriamo (2).

Se nessuna coppia di interi nella stessa riga di $\wt{T}$ condivide la medesima colonna di $T$, allora gli elementi della prima riga di $\wt{T}$ occupano colonne distinte di $T$. Esiste quindi $q_1 \in C(T)$ che li sposta tutti nella prima riga di $q_1 T$.

Iterando: per ogni $k$, si trovano $q_1, \dots, q_k \in C(T)$ tali che le entrate delle prime $k$ righe di $\wt{T}$ cadono nelle prime $k$ righe di $q_k \cdots q_1 T$. Ne segue
\[
\wt{\lambda}_1 + \cdots + \wt{\lambda}_k \leq \lambda_1 + \cdots + \lambda_k \quad \forall\, k,
\]
cioè $\lambda \geq \wt{\lambda}$ nell'ordine di dominanza. Poiché l'ordine lessicografico inverso estende quello di dominanza, $\lambda \geq^R \wt{\lambda}$. Ma per ipotesi $\lambda \not>^R \wt{\lambda}$, dunque $\lambda = \wt{\lambda}$.

Ponendo $q = q_k \cdots q_1 \in C(T)$, i tableau $qT$ e $\wt{T}$ hanno le stesse entrate in ogni riga. Esiste allora $\wt{p} \in R(\wt{T})$ con $\wt{p}\wt{T} = qT$.
}

\section{Ordine sui tableau}

\dfn{Parola di colonna}{
La \textbf{parola di colonna} $w_{\mathrm{col}}(T)$ di un tableau $T$ è la sequenza ottenuta leggendo le entrate colonna per colonna, da sinistra a destra, dall'alto verso il basso in ciascuna colonna.
}

\ex{Parola di colonna}{
  $ W_{\text{col}}(T) = 312465 $ e $ W_{\text{col}}(\wt{T}) = 264351 $
}

\dfn{Ordine sui tableau}{
Dati due tableau iniettivi $T, \wt{T}$ della stessa dimensione, si dice $\wt{T} > T$ se:
\begin{enumerate}
\item la forma di $\wt{T}$ è strettamente maggiore di quella di $T$ nell'ordine lessicografico inverso, oppure
\item le forme coincidono e il più grande intero che occupa celle diverse nei due tableau compare prima in $w_{\mathrm{col}}(\wt{T})$ che in $w_{\mathrm{col}}(T)$.
\end{enumerate}
}

\ex{Ordine su tableau}{
  Abbiamo che $ \wt{T} > T $, dato che hanno la stessa forma e il 6 compare prima in $\wt{T}$ che in $ T $.
}

\nt{
Se $T$ è un SYT, per ogni $p \in R(T)$ e $q \in C(T)$ vale $p \cdot T \geq T$ e $q \cdot T \leq T$. L'intero più grande spostato da $p$ si muove a sinistra (prima nella parola di colonna); quello spostato da $q$ scende nella colonna (dopo nella parola di colonna).
}

\cor{Ordine e intersezione righe-colonne}{
Siano $T, \wt{T}$ due SYT con $\wt{T} > T$. Allora esistono $i \neq j$ nella stessa colonna di $T$ e nella stessa riga di $\wt{T}$.
}

\pf{Dimostrazione}{
Per assurdo: se tali interi non esistono, per il Lemma di partizione esistono $\wt{p} \in R(\wt{T})$, $q \in C(T)$ con $\wt{p}\wt{T} = qT$. Poiché $T, \wt{T}$ sono SYT:
\[
T \geq q \cdot T = \wt{p} \cdot \wt{T} \geq \wt{T}
\]
cioè $T \geq \wt{T}$, in contraddizione con $\wt{T} > T$.
}

\section{Tabloidi e la rappresentazione $M^\lambda$}

\dfn{Tabloide}{
Il \textbf{tabloide} $\{T\}$ associato a un tableau iniettivo $T$ è l'orbita di $T$ sotto l'azione di $R(T)$. L'ordine degli elementi all'interno di ciascuna riga non ha importanza. L'azione di $S_n$ sui tabloidi è $\sigma \cdot \{T\} := \{\sigma \cdot T\}$.
}

\ex{Tabloide}{
  Sia $ T =  $ \begin{ytableau}
    2&4\\
    1&3&5
  \end{ytableau}, abbiamo che
  \begin{center}
  $ \{T\} =  $ \ytableausetup
{boxsize=normal,tabloids}
\ytableaushort{
24, 135
}
  \end{center}

  Inoltre, sia $ \sigma = (14) $, la sua azione naturale su $ \{T\} $ e'
  \begin{center}
    $ (14)\{T\} =  $ 
\ytableaushort{
12, 345
}
  \end{center}\ytableausetup
{boxsize=normal, notabloids}
}

\dfn{Rappresentazione $M^\lambda$}{
Lo spazio vettoriale complesso $M^\lambda$ ha come base formale l'insieme di tutti i tabloidi di forma $\lambda \vdash n$, con l'azione di $S_n$ data da permutazione dei tabloidi. Si ha l'isomorfismo
\[
M^\lambda \cong \mathrm{Ind}_{S_\lambda}^{S_n} \mathbf{1}.
\]
}

\ex{Esempi di $M^\lambda$}{
\begin{itemize}
\item $\lambda = (2,1)$, $n = 3$: i tabloidi sono $\{1,2\mid 3\}$, $\{1,3\mid 2\}$, $\{2,3\mid 1\}$, dunque $\dim M^{(2,1)} = 3$.
\item $\lambda = (n)$: un unico tabloide, $M^{(n)} \cong \mathbf{1}$ (rappresentazione banale).
\item $\lambda = (1^n)$: $n!$ tabloidi, $M^{(1^n)}$ è la rappresentazione regolare.
\end{itemize}
}

\section{Simmetrizzatori di Young}

\dfn{Simmetrizzatori}{
Sia $T$ un tableau iniettivo di dimensione $n$. Si definiscono in $\mathbb{C}[S_n]$:
\begin{enumerate}
\item \textbf{Simmetrizzatore delle righe:} $\displaystyle a_T := \sum_{p \in R(T)} p$.
\item \textbf{Simmetrizzatore delle colonne:} $\displaystyle b_T := \sum_{q \in C(T)} \mathrm{sgn}(q)\, q$.
\item \textbf{Simmetrizzatore di Young:} $c_T := b_T\, a_T$.
\end{enumerate}
}

\mprop{Proprietà algebriche dei simmetrizzatori}{
Per ogni $p \in R(T)$, $q \in C(T)$ e $\sigma \in S_n$:
\begin{enumerate}
\item $p\, a_T = a_T\, p = a_T$, \quad $a_T^2 = |R(T)|\, a_T$.
\item $q\, b_T = b_T\, q = \mathrm{sgn}(q)\, b_T$, \quad $b_T^2 = |C(T)|\, b_T$.
\item $\sigma\, a_T\, \sigma\inv = a_{\sigma \cdot T}$, \quad $\sigma\, b_T\, \sigma\inv = b_{\sigma \cdot T}$.
\end{enumerate}
}

\section{Vettore di Specht e modulo di Specht}

\dfn{Vettore di Specht}{
Il \textbf{vettore di Specht} associato a un tableau iniettivo $T$ di forma $\lambda$ è
\[
v_T := b_T \cdot \{T\} = \sum_{q \in C(T)} \mathrm{sgn}(q)\, \{q \cdot T\} \;\in M^\lambda.
\]
}

\mprop{Proprietà del vettore di Specht}{
\begin{enumerate}
\item $v_T \neq 0$ per ogni tableau iniettivo $T$.
\item Per ogni $\sigma \in S_n$: $\sigma \cdot v_T = v_{\sigma \cdot T}$.
\end{enumerate}
}

\ex{}{
  $ \lambda = (2,1) \vdash 3 $, prendiamo
  \begin{center}
    T = 
  \begin{ytableau}
  2\\
    1&3
  \end{ytableau}   
    $\wt{T} = $
  \begin{ytableau}
  3\\
    1&2
  \end{ytableau}
  \end{center}

  Calcoliamo
  \[
    a_T = id + (13)\quad b_T = -(12) + id
  \]
  \[
  a_{\wt{T}} = id + (12)\quad b_{\wt{T}} = -(13) + id
  \]
  Possiamo ora calcolare i due vettori di Specht
  \ytableausetup{boxsize=normal,tabloids}
  \begin{center}
    $ v_T = - $\ytableaushort{1,23}$ + $\ytableaushort{2,13}\\
    $ v_{\wt{T}} = - $\ytableaushort{1,23}$ + $\ytableaushort{3,12}
  \end{center}

  Che appartengono a $ M^{\lambda} $. 

}

\dfn{Modulo di Specht}{
Il \textbf{modulo di Specht} è il sottomodulo
\[
S^\lambda := \mathbb{C}[S_n] \cdot v_T \subseteq M^\lambda,
\]
generato ciclicamente da un qualsiasi vettore di Specht $v_T$ di forma $\lambda$.
}

\thm{Base standard dei moduli di Specht}{
L'insieme $\{v_T \mid T \text{ è un SYT di forma } \lambda\}$ è una base di $S^\lambda$. In particolare
\[
\dim S^\lambda = f^\lambda := \#\{\text{SYT di forma } \lambda\}.
\]
}

\ex{Modulo di Specht $S^{(2,1)}$ per $S_3$}{
Sia $T$ l'SYT $\begin{smallmatrix} 3\\ 1 & 2 \end{smallmatrix}$. Le colonne sono $\{1,3\}$ e $\{2\}$, quindi $C(T) = \{\id, (13)\}$.
\[
v_T = \{T\} - \{(13) \cdot T\} = \{1,2 \mid 3\} - \{2,3 \mid 1\}.
\]
L'altro SYT è $T' = \begin{smallmatrix} 2\\ 1 & 3 \end{smallmatrix}$ con $C(T') = \{\id, (12)\}$:
\[
v_{T'} = \{1,3 \mid 2\} - \{2,3 \mid 1\}.
\]
I vettori $v_T, v_{T'}$ sono linearmente indipendenti e formano la base standard di $S^{(2,1)}$, che è la rappresentazione standard irriducibile di $S_3$ ($\dim S^{(2,1)} = 2$), mentre $ S^{(3)} $ è la rappresentazione banale (triviale) e $ S^{(1^3)} $ è la rappresentazione segno.
}

\ex{}{
  Determinare $ S^{\lambda} $, dove $ \lambda \vdash 4 $.
}

% ============================================================
\section{Dimostrazione del Teorema della Base Standard}
% ============================================================

La dimostrazione si articola in due parti: prima mostriamo che i vettori di Specht $\{v_T \mid T \in \mathrm{SYT}(\lambda)\}$ sono linearmente indipendenti (dunque $\dim S^\lambda \geq f^\lambda$), poi con un argomento di conteggio concludiamo che $\dim S^\lambda = f^\lambda$.

\subsection{Composizioni deboli e ordine di dominanza}

\dfn{Composizione debole}{
  Una \textbf{composizione debole} (weak composition) di $n$ in $k$ parti è un vettore $\alpha = (\alpha_1, \ldots, \alpha_k) \in \mathbb{N}^k$ tale che $\alpha_1 + \cdots + \alpha_k = n$. Poniamo $\alpha_j = 0$ per $j > \ell(\alpha)$.
}

\ex{}{
  Le composizioni deboli di $3$ in $2$ parti sono: $(3,0)$, $(2,1)$, $(1,2)$, $(0,3)$.
}

\dfn{Ordine di dominanza}{
  Date due composizioni deboli $\alpha, \beta$, si dice che $\alpha \leq \beta$ nell'ordine di dominanza se
  \[
    \alpha_1 + \alpha_2 + \cdots + \alpha_i \leq \beta_1 + \beta_2 + \cdots + \beta_i \qquad \forall\, i
  \]
}

\subsection{Tabloid parziali e forma}

\dfn{Tabloid parziale}{
  Sia $\{T\}$ un tabloide di forma $\lambda \vdash n$ e sia $T$ l'unico tableau iniettivo in $\{T\}$ con righe crescenti da sinistra a destra. Il \textbf{tabloid parziale} $\{T\}^i$ è il tabloide parziale ottenuto da $T$ rimuovendo tutte le entrate $j > i$.
}

\dfn{Forma di un tabloid}{
  La \textbf{forma} $\mathrm{sh}(\{T\})$ di un tabloide $\{T\}$ è la sequenza di composizioni deboli:
  \[
    \mathrm{sh}(\{T\}) := \big(\mathrm{sh}(\{T\}^0),\, \mathrm{sh}(\{T\}^1),\, \ldots,\, \mathrm{sh}(\{T\}^n)\big)
  \]
  dove $\mathrm{sh}(\{T\}^i)$ è la composizione debole che conta quante entrate $\leq i$ compaiono in ciascuna riga.
}

\ex{}{
  Sia $\lambda = (2,2)$ e
  \begin{center}
    $T = $\begin{ytableau}
      1 & 2\\
      3 & 4
    \end{ytableau}
  \end{center}
  I tabloid parziali sono:
  \begin{center}
    $\{T\} = $\begin{ytableau}
      1 & 2\\
      3 & 4
    \end{ytableau}\quad
    $\{T\}^1 = $\begin{ytableau}
      1 & \none\\
      \none & \none
    \end{ytableau}\quad
    $\{T\}^2 = $\begin{ytableau}
      1 & 2\\
      \none & \none
    \end{ytableau}\quad
    $\{T\}^3 = $\begin{ytableau}
      1 & 2\\
      3 & \none
    \end{ytableau}\quad
    $\{T\}^4 = $\begin{ytableau}
      1 & 2\\
      3 & 4
    \end{ytableau}
  \end{center}
  Le forme (composizioni deboli) sono rispettivamente $(0,0)$, $(0,1)$, $(0,2)$, $(1,2)$, $(2,2)$.

  Quindi $\mathrm{sh}(\{T\}) = \big((0,0),\, (0,1),\, (0,2),\, (1,2),\, (2,2)\big)$.
}

\dfn{Ordine sui tabloidi}{
  Dati due tabloidi $\{T\}$ e $\{T'\}$ della stessa forma $\lambda$, si dice $\{T\} \leq \{T'\}$ se
  \[
    \mathrm{sh}(\{T\}^i) \leq \mathrm{sh}(\{T'\}^i) \qquad \forall\, i
  \]
  nell'ordine di dominanza.
}

\subsection{Lemma sull'ordine e il gruppo di colonna}

\mlenma{}{
  Sia $T$ un SYT e sia $q \in C(T)$, $q \neq \mathrm{id}$. Allora $\{T\} \not\leq \{qT\}$ nell'ordine sui tabloidi.
}

\pf{Dimostrazione}{
  Sia $q \in C(T)$ con $q \neq \mathrm{id}$. Allora $q$ sposta almeno un'entrata verso l'alto nella sua colonna (dato che $T$ è SYT, le colonne sono crescenti dal basso verso l'alto). L'entrata spostata verso l'alto va in una riga con indice maggiore, il che fa diminuire le somme parziali in almeno una componente della forma di $\{qT\}$. Quindi esiste un $i$ tale che $\mathrm{sh}(\{T\}^i) \not\leq \mathrm{sh}(\{qT\}^i)$, da cui $\{T\} \not\leq \{qT\}$.
}

\ex{Lemma sull'ordine e gruppo di colonna}{
  Sia $\lambda = (2,1)$ e consideriamo l'SYT
  \begin{center}
    $T = $\begin{ytableau}
      2\\
      1 & 3
    \end{ytableau}
  \end{center}
  Le colonne di $T$ sono $\{1,2\}$ e $\{3\}$, quindi $C(T) = \{\mathrm{id},\, (1\,2)\}$. Prendiamo $q = (1\,2) \neq \mathrm{id}$:
  \begin{center}
    $qT = $\begin{ytableau}
      1\\
      2 & 3
    \end{ytableau}
  \end{center}

  Calcoliamo le forme parziali dei due tabloidi (riga 1 = basso, riga 2 = alto):
  \[
    \begin{array}{c|c|c}
      i & \mathrm{sh}(\{T\}^i) & \mathrm{sh}(\{qT\}^i) \\
      \hline
      0 & (0,0) & (0,0) \\
      1 & (1,0) & (0,1) \\
      2 & (1,1) & (1,1) \\
      3 & (2,1) & (2,1)
    \end{array}
  \]

  Al passo $i = 1$: $\mathrm{sh}(\{T\}^1) = (1,0)$ ha prima somma parziale $1$, mentre $\mathrm{sh}(\{qT\}^1) = (0,1)$ ha prima somma parziale $0$. Poiché $1 \not\leq 0$, abbiamo $\mathrm{sh}(\{T\}^1) \not\leq \mathrm{sh}(\{qT\}^1)$ nell'ordine di dominanza, dunque $\{T\} \not\leq \{qT\}$.

  Intuitivamente: $q = (1\,2)$ ha spinto l'entrata $1$ dalla riga bassa alla riga alta, alleggerendo le somme parziali basse di $\{qT\}$ rispetto a quelle di $\{T\}$.
}

\subsection{Indipendenza lineare dei vettori di Specht}

\mprop{}{
  L'insieme $\{v_T \mid T \in \mathrm{SYT}(\lambda)\}$ è linearmente indipendente in $M^\lambda$.
}

\pf{Dimostrazione}{
  Prima osserviamo che l'ordine sui tabloidi è un \textbf{ordine totale} sull'insieme dei tabloidi di forma $\lambda$: la forma $\mathrm{sh}(\{T\})$ determina univocamente $\{T\}$ (dall'incremento $\mathrm{sh}(\{T\}^i) - \mathrm{sh}(\{T\}^{i-1})$ si ricava in quale riga si trova $i$, per ogni $i$), e la successione di composizioni $(\mathrm{sh}(\{T\}^i))_i$ è sempre comparabile fra tabloidi SYT della stessa forma.

  Supponiamo per assurdo che $\sum_{T \in \mathrm{SYT}(\lambda)} \gamma_T\, v_T = 0$ con $\gamma_T \in \mathbb{C}$ non tutti nulli. Poiché l'ordine è totale, tra tutti i tableau standard con $\gamma_T \neq 0$ esiste un \textbf{massimo} $\overline{T}$, cioè un SYT tale che $\{T\} \leq \{\overline{T}\}$ per ogni altro $T$ con $\gamma_T \neq 0$.

  Osserviamo che il tabloide $\{\overline{T}\}$ compare nell'espansione
  \[
    v_{\overline{T}} = \sum_{q \in C(\overline{T})} \mathrm{sgn}(q)\, \{q\overline{T}\}
  \]
  \textit{solo} con il termine $q = \mathrm{id}$ (coefficiente $+1$). Infatti, per ogni $q \in C(\overline{T})$ con $q \neq \mathrm{id}$, il Lemma garantisce $\{\overline{T}\} \not\leq \{q\overline{T}\}$; poiché l'ordine è totale questo significa $\{q\overline{T}\} < \{\overline{T}\}$, quindi $\{q\overline{T}\} \neq \{\overline{T}\}$.

  Affinché la somma totale faccia zero, il termine $\{\overline{T}\}$ deve essere cancellato dall'espansione di un altro vettore $v_{\widetilde{T}}$ con $\gamma_{\widetilde{T}} \neq 0$. Poiché SYT distinti producono tabloidi distinti (le righe di un SYT sono ordinate, quindi il tabloide ne determina il contenuto di ciascuna riga), $\{\overline{T}\}$ può emergere da $v_{\widetilde{T}}$ solo come termine permutato: deve valere $\{\overline{T}\} = \{\widetilde{q}\widetilde{T}\}$ per qualche $\widetilde{q} \in C(\widetilde{T})$ con $\widetilde{q} \neq \mathrm{id}$.

  Ma applicando il Lemma a $\widetilde{T}$, si ha $\{\widetilde{T}\} \not\leq \{\widetilde{q}\widetilde{T}\} = \{\overline{T}\}$; poiché l'ordine è totale, $\{\widetilde{T}\} > \{\overline{T}\}$. Questo contraddice la massimalità di $\{\overline{T}\}$: $\widetilde{T}$ ha coefficiente non nullo e il suo tabloide è strettamente maggiore di $\{\overline{T}\}$.

  Quindi nessun altro vettore può cancellare il termine $\{\overline{T}\}$: nella combinazione lineare $\sum \gamma_T v_T = 0$, il tabloide $\{\overline{T}\}$ compare con coefficiente $\gamma_{\overline{T}} \neq 0$. Ma i tabloidi sono base di $M^\lambda$, quindi linearmente indipendenti: contraddizione.
}

\subsection{Argomento di conteggio}

Per concludere la dimostrazione del teorema, serve mostrare che $\dim S^\lambda = f^\lambda$. L'indipendenza lineare da' $\dim S^\lambda \geq f^\lambda$. Il verso opposto segue da un argomento di conteggio.

\mlenma{Young}{
  Sia $\mu \vdash n$ e sia $f^\mu = \#\{\text{SYT di forma } \mu\}$. Denotiamo con $\nu \to \mu$ il fatto che $\nu$ si ottiene da $\mu$ rimuovendo una cella, e con $\lambda \leftarrow \mu$ il fatto che $\lambda$ si ottiene da $\mu$ aggiungendo una cella.
  \begin{center}
    $\nu \to \mu$:\quad
    \begin{ytableau}
      {} & {}\\
      {} & {} & {}
    \end{ytableau}
    $\;\to\;$
    \begin{ytableau}
      {} & {} & {}\\
      {} & {} & {}
    \end{ytableau}
    \qquad\qquad
    $\lambda \leftarrow \mu$:\quad
    \begin{ytableau}
      {} & {} & {}\\
      {} & {} & {}
    \end{ytableau}
    $\;\leftarrow\;$
    \begin{ytableau}
      {} & {}\\
      {} & {} & {}
    \end{ytableau}
  \end{center}
  Allora valgono le seguenti identita':
  \begin{enumerate}
    \item \textbf{Regola di discesa:} $\displaystyle\sum_{\nu \to \mu} f^\nu = f^\mu$
    \item \textbf{Regola di salita:} $\displaystyle\sum_{\lambda \leftarrow \mu} f^\lambda = (n+1)\, f^\mu$
    \item \textbf{Identita' delle dimensioni:} $\displaystyle\sum_{\lambda \vdash n} (f^\lambda)^2 = n!$
  \end{enumerate}
}

\pf{Dimostrazione}{
  \textbf{Regola di discesa (1).}
  In un qualsiasi SYT di forma $\mu$, l'elemento massimo $n$ si trova necessariamente in una cella rimovibile (un angolo esterno del diagramma). Rimuovendo la cella occupata da $n$, si ottiene un SYT valido di dimensione $n-1$ di forma $\nu \to \mu$. Questo definisce una biiezione
  \[
    \{\text{SYT di forma } \mu\} \longrightarrow \bigsqcup_{\nu \to \mu} \{\text{SYT di forma } \nu\}
  \]
  da cui $f^\mu = \sum_{\nu \to \mu} f^\nu$.

  \begin{center}
    $\mu = (2,1)$:\quad
    \begin{ytableau}
      2\\
      1 & 3
    \end{ytableau}
    $\;\longrightarrow\;$
    \begin{ytableau}
      2\\
      1 & \none
    \end{ytableau}\quad ($\nu = (1,1)$)
  \end{center}

  \textbf{Regola di salita (2).}
    Per induzione su $n$. Il caso base $n = 1$ è immediato.

  Sia $n \geq 2$ e supponiamo per ipotesi induttiva che per ogni $\nu \vdash n-1$ valga
  \[
    \sum_{\lambda \leftarrow \nu} f^\lambda = n \cdot f^\nu.
  \]
  Vogliamo mostrare che $\sum_{\lambda \leftarrow \mu} f^\lambda = (n+1) f^\mu$ per $\mu \vdash n$.

  Scriviamo $(n+1)f^\mu = f^\mu + n \cdot f^\mu$ e applichiamo la regola di discesa al primo addendo e poi al secondo:
  \[
    (n+1)f^\mu = \sum_{\nu \to \mu} f^\nu + n \sum_{\nu \to \mu} f^\nu = \sum_{\nu \to \mu} f^\nu + \sum_{\nu \to \mu} \sum_{\lambda \leftarrow \nu} f^\lambda
  \]
  dove nell'ultimo passaggio abbiamo usato l'ipotesi induttiva $n \cdot f^\nu = \sum_{\lambda \leftarrow \nu} f^\lambda$ su ogni termine.

  Lavoriamo ora sulla doppia somma. Per ogni $\nu \to \mu$, tra le partizioni $\lambda \leftarrow \nu$ ce n'è esattamente una uguale a $\mu$ (quella che si ottiene rimettendo la stessa cella che era stata tolta). Separiamo quindi i termini $\lambda = \mu$ dai termini $\lambda \neq \mu$:
  \[
    \sum_{\nu \to \mu} \sum_{\lambda \leftarrow \nu} f^\lambda = \sum_{\nu \to \mu} f^\mu + \sum_{\nu \to \mu} \sum_{\substack{\lambda \leftarrow \nu \\ \lambda \neq \mu}} f^\lambda = r \cdot f^\mu + \sum_{\nu \to \mu} \sum_{\substack{\lambda \leftarrow \nu \\ \lambda \neq \mu}} f^\lambda
  \]
  dove $r$ è il numero di celle rimovibili di $\mu$. Sostituendo:
  \[
    (n+1)f^\mu = \sum_{\nu \to \mu} f^\nu + r \cdot f^\mu + \sum_{\nu \to \mu} \sum_{\substack{\lambda \leftarrow \nu \\ \lambda \neq \mu}} f^\lambda.
  \]

  Osserviamo ora che la doppia somma conta coppie $(\nu, \lambda)$ con $\nu \to \mu$, $\lambda \leftarrow \nu$, $\lambda \neq \mu$: si toglie una cella da $\mu$ per ottenere $\nu$, poi se ne aggiunge una diversa per ottenere $\lambda \neq \mu$. Questo è equivalente ad aggiungere prima una cella a $\mu$ per ottenere $\lambda \leftarrow \mu$, poi togliere una cella diversa per ottenere $\nu \neq \mu$. Scambiando l'ordine di sommazione:
  \[
    \sum_{\nu \to \mu} \sum_{\substack{\lambda \leftarrow \nu \\ \lambda \neq \mu}} f^\lambda = \sum_{\lambda \leftarrow \mu} \sum_{\substack{\nu \to \lambda \\ \nu \neq \mu}} f^\nu.
  \]

  Applichiamo ora la regola di discesa a ogni $f^\lambda$: si ha $f^\lambda = \sum_{\nu \to \lambda} f^\nu$. Tra questi $\nu$ ce n'è esattamente uno uguale a $\mu$ (dato che $\mu \to \lambda$), quindi $\sum_{\nu \to \lambda, \nu \neq \mu} f^\nu = f^\lambda - f^\mu$. Sostituendo:
  \[
    \sum_{\lambda \leftarrow \mu} \sum_{\substack{\nu \to \lambda \\ \nu \neq \mu}} f^\nu = \sum_{\lambda \leftarrow \mu} (f^\lambda - f^\mu) = \sum_{\lambda \leftarrow \mu} f^\lambda - (r+1) f^\mu
  \]
  dove $r+1$ è il numero di celle aggiungibili a $\mu$ (un fatto geometrico: un diagramma con $r$ celle rimovibili ne ha esattamente $r+1$ aggiungibili).

  Mettendo tutto insieme:
  \[
    (n+1)f^\mu = \sum_{\nu \to \mu} f^\nu + r \cdot f^\mu + \sum_{\lambda \leftarrow \mu} f^\lambda - (r+1) f^\mu = \sum_{\lambda \leftarrow \mu} f^\lambda + \sum_{\nu \to \mu} f^\nu - f^\mu.
  \]
  Infine la regola di discesa dà $\sum_{\nu \to \mu} f^\nu = f^\mu$, quindi i due ultimi termini si cancellano e si ottiene
  \[
    (n+1)f^\mu = \sum_{\lambda \leftarrow \mu} f^\lambda
  \]

  \textbf{Identita' delle dimensioni (3).}
  Per induzione su $n$. Il caso $n = 1$ e' banale ($1^2 = 1!$). Per il passo induttivo:
  \[
    \sum_{\lambda \vdash n+1} (f^\lambda)^2 \overset{(1)}{=} \sum_{\lambda \vdash n+1} f^\lambda \left(\sum_{\nu \to \lambda} f^\nu\right)
  \]
  Invertendo l'ordine di sommazione (aggiungere a $\nu \vdash n$ anziche' togliere da $\lambda \vdash n+1$):
  \[
    = \sum_{\nu \vdash n} f^\nu \left(\sum_{\lambda \leftarrow \nu} f^\lambda\right) \overset{(2)}{=} \sum_{\nu \vdash n} f^\nu \cdot (n+1)\, f^\nu = (n+1) \sum_{\nu \vdash n} (f^\nu)^2 \overset{\text{ind.}}{=} (n+1) \cdot n! = (n+1)!
  \]
}

\pf{Dimostrazione del Teorema (Base Standard)}{
  Sappiamo che $|S_n| = n! = \sum_{\lambda \vdash n} (\dim S^\lambda)^2$ (dalla decomposizione della rappresentazione regolare e dal teorema di Maschke/Wedderburn: $\mathbb{C}[S_n] \cong \bigoplus_\lambda \mathrm{End}(S^\lambda)$, prendendo le dimensioni). Dall'indipendenza lineare si ha $\dim S^\lambda \geq f^\lambda$, quindi:
  \[
    n! = \sum_{\lambda \vdash n} (\dim S^\lambda)^2 \geq \sum_{\lambda \vdash n} (f^\lambda)^2 = n!
  \]
  La disuguaglianza e' in realta' un'uguaglianza, il che forza $\dim S^\lambda = f^\lambda$ per ogni $\lambda$. I $f^\lambda$ vettori linearmente indipendenti $\{v_T \mid T \in \mathrm{SYT}(\lambda)\}$ in uno spazio di dimensione $f^\lambda$ formano quindi una base. $\qquad\square$
}

\section{Il Lemma Tecnico}

\mlenma{Il Lemma Tecnico}{
Siano $T, \wt{T}$ tableau iniettivi di forma $\lambda, \wt{\lambda} \vdash n$, con $\lambda \not>^R \wt{\lambda}$. Allora:
\[
b_T\{\wt{T}\} = \begin{cases} 0 & \text{se } \exists\, i \neq j \text{ nella stessa colonna di } T \text{ e nella stessa riga di } \wt{T}, \\ \pm\, v_T & \text{altrimenti.} \end{cases}
\]
}

\nt{
  Il lemma tecnico implica che $b_T \cdot M^{\wt{\lambda}} = \mathbb{C}\, v_T$ per ogni $\wt{\lambda}$ con $\lambda \not>^R \wt{\lambda}$; in particolare $b_T \cdot M^\lambda = \mathbb{C}\, v_T$. Qui $\mathbb{C}\, v_T := \{z \cdot v_T \mid z \in \mathbb{C}\}$ è il sottospazio monodimensionale generato da $v_T$: il moltiplicatore $b_T$ ``proietta'' tutto $M^\lambda$ sull'asse di $v_T$.
}

\ex{}{
  $ T < \wt{T} $ e i numeri 1 e 2 sono presenti sulla stessa colonna di $ T $ e sulla stessa riga di $ \wt{T} $, allora
  \[
    b_T \{\wt{T}\} = \{3 | 12\} - \{3 | 12\} = 0
  \]
}

\pf{Dimostrazione}{
\begin{enumerate}
\item Siano $i \neq j$ nella stessa colonna di $T$ e nella stessa riga di $\wt{T}$. Sia $t = (ij)$.
\begin{itemize}
\item $t \in C(T)$, dunque $b_T t = -b_T$.
\item $t\{\wt{T}\} = \{\wt{T}\}$, perché $i,j$ stanno nella stessa riga di $\wt{T}$.
\end{itemize}
Allora:
\[
b_T\{\wt{T}\} = b_T\bigl(t\{\wt{T}\}\bigr) = (b_T t)\{\wt{T}\} = -b_T\{\wt{T}\} \implies b_T\{\wt{T}\} = 0.
\]

\item Non esistono tali interi. Per il Lemma di partizione esistono $\wt{p} \in R(\wt{T})$, $q \in C(T)$ con $\wt{p}\wt{T} = qT$. Allora:
\[
b_T\{\wt{T}\} \stackrel{\text{(i)}}{=} b_T\{\wt{p}\wt{T}\} \stackrel{\text{(ii)}}{=} b_T\{qT\} \stackrel{\text{(iii)}}{=} b_T q\{T\} \stackrel{\text{(iv)}}{=} \mathrm{sgn}(q)\,b_T\{T\} \stackrel{\text{(v)}}{=} \mathrm{sgn}(q)\,v_T = \pm v_T
\]
dove: (i) $\wt{p} \in R(\wt{T})$ non altera il tabloide; (ii) si usa $\wt{p}\wt{T} = qT$; (iii) $q$ agisce sul tableau dentro il tabloide; (iv) $q \in C(T) \implies b_T q = \mathrm{sgn}(q)\,b_T$; (v) $v_T = b_T\{T\}$ per definizione.
\end{enumerate}
}

\cor{Annullamento per tableau standard}{
Se $T, \wt{T}$ sono SYT con $\wt{T} > T$, allora $b_T\{\wt{T}\} = 0$.
}

\pf{Dimostrazione}{
Dall'ipotesi $\wt{T} > T$ (nell'ordine sui tableau iniettivi della stessa forma) e dal fatto che entrambi sono SYT della stessa forma $\lambda$, il \textbf{Lemma di partizione dei tableau} garantisce l'alternativa: o esistono $i \neq j$ nella stessa colonna di $T$ e nella stessa riga di $\wt{T}$, oppure esiste $\wt{p} \in R(\wt{T})$, $q \in C(T)$ con $\wt{p}\wt{T} = qT$. Quest'ultima opzione è impossibile: se $\wt{p}\wt{T} = qT$ allora $\wt{T}$ e $T$ avrebbero le stesse entrate per righe (a meno di ordine), ma SYT distinti della stessa forma producono tabloidi distinti (e quindi row-set distinti per qualche riga), contraddizione. Dunque si è nel primo caso, e si applica il Caso 1 del Lemma Tecnico.
}

\section{Classificazione delle rappresentazioni irriducibili di $S_n$}

\thm{Classificazione delle rappresentazioni irriducibili di $S_n$}{
I moduli di Specht $\{S^\lambda \mid \lambda \in \mathrm{Par}(n)\}$ costituiscono un sistema completo di rappresentazioni irriducibili di $S_n$, a due a due non isomorfe.
}

\pf{Dimostrazione}{
\begin{enumerate}
\item \textbf{Non banalità.} Il vettore $v_T \neq 0$: nell'espansione $v_T = \sum_{q \in C(T)} \mathrm{sgn}(q)\{qT\}$, il termine $q = \id$ produce $\{T\}$ con coefficiente $+1$, e nessun altro termine lo cancella. Dunque $S^\lambda \neq 0$.

\item \textbf{Non isomorfismo.} Sia $T$ un tableau di forma $\lambda$. L'azione di $b_T$ su $S^\lambda$ soddisfa
\[
b_T \cdot S^\lambda = \mathbb{C}\, v_T \neq 0.
\]
Se $\wt{\lambda} >_{\mathrm{lex}} \lambda$, il Lemma Tecnico (Caso 1) implica $b_T \cdot S^{\wt{\lambda}} = 0$. Questa differenza nell'azione algebrica esclude ogni isomorfismo tra $S^\lambda$ e $S^{\wt{\lambda}}$.

\item \textbf{Irriducibilità.} Per assurdo, sia $S^\lambda = V \oplus W$. Applicando $b_T$:
\[
\mathbb{C}\, v_T = b_T V \oplus b_T W.
\]
Essendo $\mathbb{C}\, v_T$ unidimensionale, uno dei due sommandi è nullo; supponiamo $b_T W = 0$. Allora $v_T \in b_T V$, ma $v_T$ genera ciclicamente $S^\lambda$, dunque $V = S^\lambda$.

\item \textbf{Completezza.} La famiglia $\{S^\lambda\}_{\lambda \vdash n}$ è indicizzata dalle partizioni di $n$, che sono in biiezione con le classi di coniugio di $S_n$. Poiché il numero di rappresentazioni irriducibili di un gruppo finito coincide con il numero delle sue classi di coniugio, la famiglia è completa.
\end{enumerate}
}

\subsection{Ruolo nella teoria}

La formula-chiave che governa l'intera classificazione è l'asimmetria dell'azione di $b_T$:
\[
b_T \cdot S^\lambda = \mathbb{C}\, v_T \neq 0, \qquad b_T \cdot S^{\wt{\lambda}} = 0 \;\;\text{se } \wt{\lambda} >^R \lambda.
\]
Questa proprietà garantisce simultaneamente la non isomorfismo tra moduli associati a partizioni diverse e l'impossibilità di scomporre ciascun modulo in sottomoduli propri.
